\qquad During the internal penetration test against \enterpriseName, \brandName identified  \the\numexpr\totalCount\relax findings that
threaten the confidentiality, integrity, and availability of \enterpriseName’s information systems. The findings were
categorized by severity level, with \countCritical of the findings being assigned a critical-risk rating,\countHigh high-risk, \countMedium medium-risk, 
\countLow low risk, and \countInfo info risk. There was also one (1) informational finding related to enhancing security monitoring capabilities within the
internal network.
The tester found \enterpriseName’s patch and vulnerability management to be well-maintained. None of the findings in this
report were related to missing operating system or third-party patches of known vulnerabilities in services and
applications that could result in unauthorized access and system compromise. Each flaw discovered during testing was
related to a misconfiguration or lack of hardening, with most falling under the categories of weak authentication and
weak authorization.
One finding involved a network communication protocol that can be “spoofed” to retrieve passwords for internal users
that can be used to gain unauthorized access if an attacker can gain unauthorized access to the network without
credentials. In most corporate environments, this protocol is unnecessary and can be disabled. It is enabled by default
primarily for small and medium sized businesses that do not have the resources for a dedicated hostname resolution (the
“phonebook” of your network) server. During the assessment, the presence of these resources was observed on the
network, so \enterpriseName should begin formulating a test plan to disable the dangerous service.
The next issue was a weak configuration involving service accounts that allows any authenticated user to steal a
component of the authentication process that can often be guessed offline (via password “cracking”) to reveal the
human-readable form of the account’s password. These types of service accounts typically have more privileges than a
standard user, so obtaining one of their passwords in clear text could result in lateral movement or privilege escalation
and eventually in complete internal network compromise. The tester also noticed that the same password was used for
administrator access to all servers within the internal network. This means that if one server is compromised, an attacker
can re-use this password to access any server that shares it for administrative access. Fortunately, both issues can be
corrected without the need for third-party tools. Microsoft’s Active Directory contains settings that can be used to
minimize the risk of these resources being abused for the benefit of malicious users.
A webserver was also found to be running a web application that used weak and easily guessable credentials to access an
administrative console that can be leveraged to gain unauthorized access to the underlying server. This could be
exploited by an attacker on the internal network without needing a valid user account. This attack is very welldocumented, so it is an exceedingly likely target can be particularly damaging, even in the hands of an unskilled attacker.
Ideally, direct external access to this service would be disabled, but if it cannot be, it should be reconfigured with
exceptionally strong credentials that are rotated frequently. \enterpriseName may also want to consider maximizing the log
data collected from this device to ensure that attacks against it can be detected and triaged quickly.
7
The tester also found shared folders with excessive permissions, meaning that all users in the internal network can access
a considerable amount of data. While sharing files internally between departments and users is important to day-to-day
business operations, wide open permissions on file shares may result in unintentional disclosure of confidential
information. Even if a file share does not contain any sensitive information today, someone may unwittingly put such
data there thinking it is protected when it isn’t. This configuration should be changed to ensure that users can access only
what is necessary to perform their day-to-day duties.
Finally, the tester noticed that testing activities seemed to go mostly unnoticed, which may represent an opportunity to
improve visibility into the internal network and indicates that a real-world attacker might remain undetected if internal
access is achieved. \enterpriseName should create a remediation plan based on the Remediation Summary section of this
report, addressing all high findings as soon as possible according to the needs of the business. \enterpriseName should also
consider performing periodic vulnerability assessments if they are not already being performed. Once the issues
identified in this report have been addressed, a more collaborative, in-depth Active Directory security assessment may
help identify additional opportunities to harden the Active Directory environment, making it more difficult for attackers
to move around the network and increasing the likelihood that \enterpriseName will be able to detect and respond to
suspicious activity